\section{Scope, attribution and verification}
All four principal results remain attributed to Hou--Feng--Leemans. The earlier project proof supplied the explicit second affine model, the common period-lattice argument, exact presentation identifications of the extensions, and an independently reproduced nonexistence base. This edition incorporates those proofs and their full dependencies; it does not claim those steps again as new results. The source's cyclic-core descent is retained with its hypotheses and quotient step made explicit. The mixed-prime type argument is a short reworking of its solvable-group argument. Neither reuse nor a shorter presentation establishes historical priority.

The common paper [U] includes the complete supported-extension and square proofs. The two cubic kernel computations and the order-384 catalogue search are exact finite computations with declared external algorithms. The three previously verified Lean declarations prove the literal sign obstruction; the previously introduced two-declaration affine component proves the preservation of an intersection under all smaller normal periods. Four reused generic quotient criteria are included as its formal dependencies. Affine coordinates, actual subgroup presentations, finite catalogue completeness and the full geometric classification remain outside these Lean statements. Every mathematical program is supplied in full. The complete common paper, application and required programs must be charged in a later comparison; no factor is inferred from page or title counts.

The source's broader Problem 1.1, the characterization of all groups of
order \(2^n\) or \(2^n p\) arising from regular or chiral polytopes, is not
solved here. Its final question about type (2) for primes at least five
has already been answered in the literature, as recorded in the
project's question audit and the explicit equal-label construction in Section 9, with its attribution to [MW20] and [LZ25]. That existing answer earns no new question credit here. Introductory
survey statements, the small-order table as an independently complete
census, and ancillary uniqueness remarks about the two order384 seeds
are outside this principal contract.
